% Generated draft: requires mathematical and source review.
\documentclass[11pt]{article}
\usepackage[margin=25mm]{geometry}
\usepackage{fontspec}
\setmainfont{Latin Modern Roman}
\usepackage{amsmath,amssymb,mathtools,unicode-math}
\setmathfont{Latin Modern Math}
\usepackage{enumitem,needspace,xurl,hyperref}
\setlength{\parindent}{0pt}
\setlength{\parskip}{5pt}
\setlength{\emergencystretch}{3em}
\newcommand{\N}{\mathbb{N}}
\newcommand{\Z}{\mathbb{Z}}
\newcommand{\Q}{\mathbb{Q}}
\newcommand{\R}{\mathbb{R}}
\newcommand{\C}{\mathbb{C}}
\newcommand{\F}{\mathbb{F}}
\newcommand{\A}{\mathbb{A}}
\newcommand{\PP}{\mathbb{P}}
\newcommand{\E}{\mathbb{E}}
\newcommand{\eps}{\varepsilon}
\newcommand{\dd}{\,\mathrm{d}}
\begin{document}
\section*{UnsolvedMath ID: 30003696}
\textbf{Source problem code:} catalogue entry
\subsection{Definitions and mathematical statement}
The supplied record does not explicitly specify the ground field or parameter space. The following gives a precise formulation under the conventional interpretation of smooth complex plane cubics parametrized by their equations; this interpretation requires confirmation against the source.

Let \(V=\operatorname{Sym}^{3}((\C^{3})^{*})\), and give \(\PP(V)\) its complex analytic topology. Define
\[
B=\{[f]\in\PP(V): C_f=\{x\in\PP^{2}(\C):f(x)=0\}\text{ is smooth}\}.
\]
Thus equations differing by a nonzero scalar determine the same point of \(B\). The universal family is
\[
\mathcal C=\{([f],x)\in B\times\PP^{2}(\C):f(x)=0\},
\qquad \pi:\mathcal C\longrightarrow B,\quad ([f],x)\longmapsto[f].
\]
For each integer \(n\geq1\), let
\[
\mathcal F_n=\{([f],x_1,\ldots,x_n)\in B\times(\PP^{2}(\C))^n:
 f(x_i)=0\text{ for all }i,\ x_i\ne x_j\text{ for }i\ne j\}.
\]
Give \(\mathcal F_n\) the subspace topology and let
\[
\mathcal U_n=\mathcal F_n/\mathfrak S_n,
\qquad q_n:\mathcal U_n\longrightarrow B,
\]
where the symmetric group permutes the points and \(q_n\) remembers the cubic. A continuous multisection of degree \(n\) is a continuous section \(s:B\to\mathcal U_n\) of \(q_n\). Its value assigns an unordered set of exactly \(n\) distinct points of \(C_f\) to each \([f]\in B\).

Two such sections \(s_0,s_1\) are homotopic through multisections if there is a continuous map \(H:B\times[0,1]\to\mathcal U_n\) satisfying
\[
H(b,0)=s_0(b),\qquad H(b,1)=s_1(b),\qquad q_n(H(b,t))=b
\quad(b\in B,\ t\in[0,1]).
\]
In particular, the points must remain distinct throughout the homotopy.

\textbf{Classification problem.} Determine all integers \(n\geq1\) for which \(q_n\) admits a continuous section, and for every such \(n\), classify its sections up to homotopy through sections.

This formulation uses embedded complex cubics with fixed ambient projective plane. It imposes no additional equivariance under projective coordinate changes. The supplied record does not explicitly resolve whether this parameter space, a quotient moduli space, or a moduli stack is intended; these alternatives must not be identified without further source verification.
\subsection{Short English statement}
For each positive integer \(n\), determine whether one can continuously assign an unordered set of \(n\) distinct points to every smooth plane cubic, and classify all assignments up to continuous deformation that preserves distinctness. The displayed formulation assumes complex cubics parametrized by equations; the supplied record leaves that choice of domain implicit.
\subsection{Sources}
\begin{itemize}
\item Local UnsolvedMath registry: supplied complete record for permanent ID \(30003696\), titled \emph{Continuous Choices of Points on Smooth Cubic Curves}, source problem code \texttt{catalogue entry}, in the Oberwolfach Reports Open Problems collection. The record gives proposed year \(2018\) and names no proposer.
\item Supplied report citation: \emph{Topology of Arrangements and Representation Stability (2019)}, Oberwolfach Reports, DOI \texttt{10.4171/owr/2018/2}, pp.\ 43--123. \url{https://doi.org/10.4171/owr/2018/2}. The background describes Oberwolfach Reports as documenting mathematical workshops and their open-problem discussions.
\item Supplied literature reference: Banerjee and Chen, cubic plane curves. \url{https://arxiv.org/abs/2101.03824}.
\item Supplied literature reference: Yao, rigidity of cubic multisections. \url{https://arxiv.org/abs/2607.14128}.
\item Supplied literature reference: Oberwolfach report. \url{https://doi.org/10.4171/owr/2018/2}.
\item Literature assessment dated \texttt{2026-08-21}: the supplied curation labels the problem \emph{partially solved} and reports nonexistence for many degrees, uniqueness at degree nine, nonexistence at degree eighteen, and newer rigidity at degree twenty-seven, while stating that classification for every positive degree remains open. This is supplied dated curation, not an independently verified literature assessment. The background explicitly requests independent verification of time-sensitive status and cited claims.
\item The supplied registry labels the statement \texttt{corrected\_verified} and says that the official report was checked and this is its canonical multisection-classification formulation. These are supplied provenance claims; no source visit or independent mathematical verification is asserted here.
\end{itemize}
\end{document}
